\chapter{IMPLEMENTATION}

\justify

This chapter describes the implementation methodology and key modules of the HPDC Defect Detection System. The implementation is divided into five functional modules: dataset analysis, patch generation, dataset engineering, RF-DETR training, and evaluation and inference.

\section {Proposed Methodology (Explain the methodology with diagrams/
Algorithms/flowchart explanation)}
\justify

The core methodology is a \textbf{patch-based sliding-window training pipeline} designed for extreme small-object detection. Since crack bounding boxes occupy less than 0.1\% of the full image area, full-image detection is infeasible. Each HPDC image is divided into overlapping 1024$\times$1024 patches using a stride of 665 px (35\% overlap), with Shapely polygon-aware annotation clipping to preserve accurate ground-truth boundaries.

\begin{figure}[H]
\centering
\begin{tikzpicture}[node distance=0.5cm,
block/.style={rectangle, rounded corners=4pt, draw=magnaDark,
fill=magnaLightGray, text width=7.5cm,
minimum height=1cm, align=center, font=\small},
arrow/.style={-Stealth, thick, color=magnaRed}]
\node[block, fill=magnaBlue!15, draw=magnaBlue] (input)
    {\textbf{Input:} Full image + COCO polygon annotations};
\node[block, below=0.5cm of input] (stride)
    {\textbf{Sliding Window:} Stride $= 1024 \times (1 - 0.35) = 665$\,px};
\node[block, below=0.5cm of stride] (extract)
    {\textbf{Patch Extraction:} Crop \texttt{image[y:y+H, x:x+W]}};
\node[block, below=0.5cm of extract] (shapely)
    {\textbf{Polygon Clipping (Shapely):} Intersect polygon with patch boundary;
    recompute bbox from clipped polygon};
\node[block, below=0.5cm of shapely] (balance)
    {\textbf{Balancing:} Downsample negatives to 1:2 pos:neg ratio};
\node[block, below=0.5cm of balance] (aug)
    {\textbf{Augmentation:} Horizontal flip on crack patches only};
\node[block, fill=successGreen!10, draw=successGreen, below=0.5cm of aug] (output)
    {\textbf{Output:} COCO-format JSON + PNG patches for training};
\draw[arrow] (input)   -- (stride);
\draw[arrow] (stride)  -- (extract);
\draw[arrow] (extract) -- (shapely);
\draw[arrow] (shapely) -- (balance);
\draw[arrow] (balance) -- (aug);
\draw[arrow] (aug)     -- (output);
\end{tikzpicture}
\caption{Implementation methodology flowchart: patch-based data engineering pipeline.}
\label{fig:methodology}
\end{figure}

The EMA training strategy smooths weight updates:
\begin{equation}
\boldsymbol{\theta}_{\text{EMA}} = \alpha \cdot \boldsymbol{\theta}_{\text{EMA}} + (1-\alpha) \cdot \boldsymbol{\theta}_{\text{current}}
\end{equation}
where $\alpha \approx 0.9999$. EMA weights consistently outperformed instantaneous weights across all experiments.

\section {Description of Modules (Briefly describe  important modules
Modules descriptions
Module Name :  module( )
Input :
Output :
 Describe Code/ Pseudocode/ flowchart
}

\subsection{Module 1: Dataset Analysis Module}

\textbf{Module Name:} \texttt{dataset\_analysis()}

\textbf{Input:} Raw COCO JSON annotation file for the full 256-image dataset.

\textbf{Output:} Statistical summary of crack/cold shut sizes, class distribution, annotation quality report.

\textbf{Description:} This module performs exploratory data analysis (EDA) to characterise the defect dataset. It computes bounding box area distributions, identifies that crack areas occupy a median of 0.065\% of image area, and validates polygon annotation quality. The EDA confirmed that full-image training is infeasible and a patch-based strategy is required.

\begin{lstlisting}[language=Python, caption={Dataset analysis pseudocode}]
def dataset_analysis(coco_json_path):
    coco = load_coco(coco_json_path)
    image_area = 4096 * 3000  # pixels^2
    for ann in coco['annotations']:
        bbox_area = ann['bbox'][2] * ann['bbox'][3]
        relative_area = bbox_area / image_area * 100
        record(ann['category_id'], relative_area)
    # Output: median crack area ~ 0.065%, Cold Shut ~ 0.12%
    plot_distributions()
\end{lstlisting}

\subsection{Module 2: Patch Generation Module}

\textbf{Module Name:} \texttt{generate\_patches()}

\textbf{Input:} Full HPDC image (4096$\times$3000 PNG), COCO annotation JSON.

\textbf{Output:} Set of 1024$\times$1024 PNG patches + updated COCO JSON with clipped annotations.

\textbf{Description:} Enumerates sliding-window positions with stride 665 px. For each patch, clips all polygon annotations using Shapely geometry intersection and recomputes bounding boxes. Discards annotations where intersection area falls below a minimum threshold.

\begin{lstlisting}[language=Python, caption={Patch generation pseudocode}]
def generate_patches(image, annotations, patch_size=1024, overlap=0.35):
    stride = int(patch_size * (1 - overlap))  # 665 px
    for y in range(0, H - patch_size + 1, stride):
        for x in range(0, W - patch_size + 1, stride):
            patch = image.crop(x, y, patch_size, patch_size)
            patch_box = Polygon([(x,y),(x+S,y),(x+S,y+S),(x,y+S)])
            clipped = []
            for ann in annotations:
                poly  = Polygon(ann['segmentation'])
                inter = poly.intersection(patch_box)
                if inter.area > MIN_THRESHOLD:
                    bbox = inter.bounds  # recomputed bbox
                    clipped.append(make_coco_ann(bbox, ann['category_id']))
            save_patch(patch, clipped)
\end{lstlisting}

\subsection{Module 3: Dataset Engineering Module}

\textbf{Module Name:} \texttt{balance\_and\_augment()}

\textbf{Input:} All generated patches with annotation counts.

\textbf{Output:} Balanced and augmented patch dataset ready for training.

\textbf{Description:} Separates positive (defect-containing) and negative (background) patches. Randomly downsamples negatives to a 1:2 positive:negative ratio. Applies horizontal flip to crack-containing patches, recomputing bounding box coordinates as $x_{\text{new}} = W - x - w$.

\begin{lstlisting}[language=Python, caption={Balancing and augmentation pseudocode}]
def balance_and_augment(patches):
    pos = [p for p in patches if p.annotations]
    neg = [p for p in patches if not p.annotations]
    neg = random.sample(neg, min(len(neg), 2 * len(pos)))
    augmented = []
    for patch in pos:
        if has_crack(patch):
            flipped_img  = cv2.flip(patch.image, 1)
            flipped_anns = [flip_bbox(a, patch.width) for a in patch.annotations]
            augmented.append(Patch(flipped_img, flipped_anns))
    return pos + neg + augmented

def flip_bbox(ann, W):
    x, y, w, h = ann['bbox']
    return {'bbox': [W - x - w, y, w, h], 'category_id': ann['category_id']}
\end{lstlisting}

\subsection{Module 4: RF-DETR Training Module}

\textbf{Module Name:} \texttt{train\_rfdetr()}

\textbf{Input:} COCO-format patch dataset (train/val splits), training configuration.

\textbf{Output:} Best EMA checkpoint (\texttt{checkpoint\_best\_ema.pth}), per-epoch mAP metrics.

\textbf{Description:} Fine-tunes RF-DETR Medium with DINOv2 ViT-B backbone using Hungarian bipartite matching loss (classification + $\mathcal{L}_1$ + GIoU). EMA weights are maintained and saved when validation EMA mAP50-95 improves. Best result achieved: EMA mAP50-95 = 0.2379 at Epoch 10.

\begin{lstlisting}[language=Python, caption={RF-DETR training configuration (Experiment 7)}]
model = RFDETRMedium(pretrain_weights='rf-detr-medium.pth')
model.train(
    dataset_dir = 'RF_DETR_Medium_1024/',
    epochs      = 20,
    batch_size  = 12,
    imgsz       = 1024,
    workers     = 16,
    device      = 'cuda'
)
# Results: EMA mAP50-95 = 0.2379, mAP50 = 0.4647 (Epoch 10)
\end{lstlisting}

\subsection{Module 5: Evaluation and Inference Module}

\textbf{Module Name:} \texttt{evaluate()} / \texttt{detect()}

\textbf{Input:} Test split patches (COCO format) or full HPDC image + trained EMA checkpoint.

\textbf{Output:} COCO mAP50-95, mAP50, per-class AP (evaluation) or annotated image with bounding boxes (inference).

\textbf{Description:} Evaluation runs the EMA checkpoint against the test split using pycocotools for standard COCO metrics. Inference generates patches from a full input image, runs RF-DETR forward pass on each patch, maps bounding boxes back to full-image coordinates, and renders the annotated output.

\begin{lstlisting}[language=Python, caption={Inference pseudocode}]
def detect(image_path, checkpoint, conf_threshold=0.3):
    model = RFDETRMedium()
    model.load_state_dict(torch.load(checkpoint))
    model.eval()
    image   = load_image(image_path)
    patches = generate_patches(image, patch_size=1024, overlap=0.35)
    results = []
    for (patch, x_off, y_off) in patches:
        dets = model(patch)
        for det in dets:
            if det.confidence > conf_threshold:
                det.bbox = offset_bbox(det.bbox, x_off, y_off)
                results.append(det)
    return annotate_image(image, results)
\end{lstlisting}
